\documentclass[10pt,a4paper]{article}
\usepackage[T1]{fontenc}
\usepackage[utf8]{inputenc}
\usepackage{lmodern,amsmath,amssymb,array,longtable,booktabs,enumitem}
\usepackage{kotex}
\usepackage[margin=20mm]{geometry}
\usepackage[unicode,colorlinks=true,urlcolor=blue,linkcolor=blue]{hyperref}
\setlength{\parindent}{0pt}
\setlength{\parskip}{5pt}
\setlength{\emergencystretch}{2em}
\newcolumntype{P}[1]{>{\raggedright\arraybackslash}p{#1}}
\title{기본 모델: 파라미터 근거 및 설정 이유}

\begin{document}
\maketitle


\section{개요}
baseline model에 대한 합성 데이터 기반 모델 검증과 민감도 실험을 뒷받침하는 문서다. \textbf{기본 모델은 고정된 $M_c$, \textcolor{red}{외생적으로 주어지는 $p_c^s$}, 도시 진입 결정 $y_c$, 정수 차량 배치량 $q_c$, 그리고 기대 초과 용량 및 미충족 수요 비용으로 구성}된다.

\textbf{근거 구분:} %표로 정리%
O는 해당 출처의 특정 표본에서 보고된 관측값, E는 출처에서 발표한 추정값, D는 산술적으로 환산한 값, A는 연구자가 설계한 합성 실험의 가정을 뜻한다. 

%표 아래 문구%
\textcolor{blue}{※ 규모를 참고할 만한 출처가 있다는 이유만으로 제안 범위가 실증적으로 확인된 범위가 되는 것은 아니다. 단지 참고만 하고 임의로 산출하여 범위를 계산하였다.} \\

\textbf{평가 기준:} \textbf{도시 진입 여부와 차량 대수를 한 번 결정하고 1년간 운영하는 상황을 평가한다.} \(F_c\)는 진입에 따른 일회성 준비·인프라 구축비(USD), \(K_c\)는 차량당 연간 추가 고정비(USD/대·년)이다. \(c_c^+\)는 미사용 용량에 대한 금전 환산 페널티, \(c_c^-\)는 미충족 통행의 상실 공헌이익이며, 모두 USD/건으로 설정한다. 수요·처리량은 연간 건수를 사용하며, 목적함수는 초기비용과 1년간의 비용·손실·페널티를 합산한다.

\section{파라미터 범위와 권장 설정값}
\small
\begin{longtable}{P{0.10\linewidth}P{0.34\linewidth}P{0.48\linewidth}}
\toprule 기호 & 합성 실험 설정 / 정의역 & 근거와 해석 \\
\midrule\endfirsthead
\toprule 기호 & 합성 실험 설정 / 정의역 & 근거와 해석 \\
\midrule\endhead

$C$ & \textbf{도시 A, B, C의 3개} & \textbf{A. 임의 가정}. 모든 후보 도시에서 운영이 가능하다고 가정 \\
\midrule

$S$ & \textbf{시나리오 3개} 이후 여러 난수 시드로 50, 100, 500개의 시나리오를 표본 추출. & \textbf{A. 임의 가정.} \\ \midrule

$1/|S|$ & 모든 시나리오의 발생 확률을 동일 가정. & \textbf{A. 임의 가정.} \\ \midrule

$y_c$ & \textbf{$\{0,1\}$}. & \textbf{최적화로 결정.} $F_C > 0$이면 최적해에서 $q_c=0$인 도시는 $y_c=0$이다. \\ \midrule

$q_c$ & \textbf{$0 - \min(\bar q_c,Q)$} 사이 정수 ($0\le q_c\le\min(\bar q_c,Q)$). &  \textbf{A. 임의 가정.} \\ \midrule

$\bar q_c$ & 도시당 500대. (비교 후보: 100, 500, 1000대) & \textbf{A/O. 임의가정/관측된 Waymo 공식 발표 참고}. \hyperlink{rationale-R1}{[R1]}. \\ \midrule

$Q$ & 총 1000대. (비교 후보: 200, 1000, 2000대.) & \textbf{A. 임의 가정 (기업의 자원 제약).} \textcolor{red}{단, $M_c\bar p_c/\kappa_c$로 필요한 차량 규모를 점검해야 한다.} \\ \midrule

$\kappa_c$
& 차량 1대당 연간 $8{,}760(\text{건}/(\text{대}\cdot\text{년}))$.
(비교 후보: $\{6{,}570,\ 8{,}760,\ 10{,}950\}\, \text{건}/(\text{대}\cdot\text{년})$).
& \textbf{D/O. 산출/관측된 Waymo 공식 발표 참고.}
민감도 분석에서는 기본값의 $\pm25\%$ 적용. \hyperlink{rationale-R2}{[R2]}. \\ \midrule


$M_c$ & 연간 $770{,}544{,}140(\text{건}/\text{년})$의 잠재 이용 기회. (비교 후보: 577,908,105, 770,544,140, 963,180,175). & \textbf{D/O. 산출/관측된 자료 참고 (SFCTA의 연간 Uber/Lyft 집계, SFMTA 2019년 통행분담률 조사.} 두 자료를 활용하여 실험용 기본값을 산출하고 민감도 분석에서는 기본값의 $\pm25\%$ 적용. \hyperlink{rationale-R3}{[R3]}, \hyperlink{rationale-R4}{[R4]} \\ \midrule

$p_c^s$ & 정의역 $[0,1]$. A: 탐색 범위 0.02--0.20. 기본 평균값 $(0.10,0.08,0.06)$. & \textbf{A. 임의가정.} 원래 추정하는 값이지만 모델의 검증을 위해 임의 가정. \\ \midrule

$F_c$ & \$3,500,000.  비교 후보: \{\$2,625,000,\ \$3,500,000,\ \$4,375000\}\ & \textbf{D/O. 산출/관측된 자료 (Waymo Austin phase)} 기존 건물을 임차해 운영 거점 하나를 개조하는 진입 방식을 가정하여 충전/전기 공사 + 운영거점 개조 + 추가 준비비용($\pm25\%$ 임의 가정)으로 기본값 산출 \hyperlink{rationale-R5}{[R5]} \\ \midrule

$K_c$ & \$7972. 비교 후보: \{\$5979,\ \$7972,\ \$9965\}\  & \textbf{D/O. 산출/관측된 자료 (AAA).} 개인용 신차 보유 비용(감가상각비, 금융비용, 보험료, 면허/등록/세금) 등의 연간 합계 참고하여 기본값 산출 \hyperlink{rationale-R6}{[R6]} \\ \midrule

$c_c^+$ & 기본값은 미사용 통행 용량 1건당 \$0. 선택적인 패널티 수준은 1, 3, 5. & \textbf{A.임의 가정.} 별도의 추가 유휴 용량 손실을 가정하지 않는다. \textcolor{red}{양수면 회사의 선호를 반영한 패널티를 부여한 것이다.} \\ \midrule

$c_c^-$ & 미충족 통행 1건당 기본값 \$10. (비교 후보: 2, 5, 10, 15, 20). & \textbf{D/O. 관측된 자료 참고 (Obi: The road ahead} 공헌이익은 기업 귀속 운임에서 통행 발생에 따른 추가 비용을 뺀 값이다. \hyperlink{rationale-R7}{[R7]} \\ \midrule

\bottomrule
\end{longtable}
\normalsize

\section{출처별 원래 수치와 적용 한계}
\subsection*{\hypertarget{rationale-R1}{R1. 차량 배치 규모: Uber 및 Waymo 공개 발표}}
\textbf{O:} Uber는 2025년 6월 24일 Austin의 Uber 플랫폼에서 Waymo 차량 100대가 운영 중이라고 밝혔다. Waymo는 2025년 5월 5일 San Francisco, Los Angeles, Phoenix, Austin의 4개 도시에서 상업용 차량 1,500대 초과를 운영한다고 밝혔고, 2026년 2월 17일에는 회사 전체 차량 규모를 3,000대라고 밝혔다. 
\\ \textbf{적용:} 관측 수치는 차량 규모를 비교하는 참고값이다. 연구자 가정(A)으로 도시당 상한 $\bar q_c=500$대와 총 가용량 $Q=1{,}000$대를 기본값으로 두고, 각각 $\{100,500,1{,}000\}$대와 $\{200,1{,}000,2{,}000\}$대를 비교한다.
\\ \textbf{한계:} Austin의 100대는 Uber 플랫폼의 당시 수치이며 도시 전체 차량 수가 아니다. 1,500대 초과는 4개 도시의 합계이고, 3,000대는 2026년 2월 회사 전체 차량 규모로 집계 범위와 시점이 다르다. 어느 자료도 도시별 물리적 상한이나 본 실험의 최적 배치량을 제시하지 않는다. 또한 Waymo 한 사업자의 차량 수를 로보택시 전체 공급량으로 해석할 수 없다.
\par\href{https://investor.uber.com/news-events/news/press-release-details/2025/Atlanta-The-Future-is-Here-with-Waymo-and-Uber-2025-y1tuDCCSgu/default.aspx}{Uber: Atlanta, The Future is Here with Waymo and Uber.} \par\href{https://waymo.com/blog/2025/05/scaling-our-fleet-through-us-manufacturing/#:~:text=over\%201\%2C500\%20vehicles}{Waymo Team, \textit{Scaling our fleet through U.S. manufacturing}, 2025년 5월 5일, Tooling up for growth 첫 문단(4개 도시의 상업용 차량).} \par\href{https://waymo.com/blog/shorts/advice-not-control-the-role-of-remote-assistance/#:~:text=fleet\%20of\%203\%2C000\%20vehicles}{Waymo, \textit{Advice, not control: the role of Remote Assistance in Waymo's operations}, 2026년 2월 17일, 마지막 본문 문단(전체 차량 규모).}

\subsection*{\hypertarget{rationale-R2}{R2. 차량당 연간 처리량: Waymo 운영 실적 참고}}

\noindent
\textbf{O:}
Waymo는 Phoenix, San Francisco, Los Angeles, Austin의
4개 도시에서 주간 유료 운행 250,000건 초과와
상업용 차량 1,500대 초과를 보고하였다. \\
\textbf{적용:}
미국 로보택시의 차량당 운행 규모를 참고하여
$\kappa_c$의 실험용 기본값을 설정한다. \\
\textbf{한계:}
2025년 5월 보고서 자료이며 두 수치 모두 `초과'로 제시되어 정확한 차량당 평균을 산출할 수 없으며, 실제 운행 실적은 최대 처리용량과 다르다. 도시별 차이와 연중 변동을 반영한 수치도 아니다.

\medskip
\noindent
\href{https://waymo.com/blog/2025/05/scaling-our-fleet-through-us-manufacturing/#:~:text=more\%20than\%20250\%2C000\%20paid\%20trips\%20each\%20week}
{Waymo Team, \textit{Scaling our fleet through U.S. manufacturing},
2025년 5월 5일, 본문 첫 문단(주간 운행 건수)};
\href{https://waymo.com/blog/2025/05/scaling-our-fleet-through-us-manufacturing/#:~:text=over\%201\%2C500\%20vehicles}
{같은 자료, ``Tooling up for growth'' 첫 문단(상업용 차량 수)}.

\medskip
\noindent
\textbf{D: 산출 및 실험 설정.}
공개 총량의 단순 비율
$250{,}000/(1{,}500\times7)\approx23.8$을 참고하여
기본값을 하루 24건/대로 정한다.
운행 건수와 차량 수 모두 같은 4개 도시의 합계이므로
4로 추가로 나누지 않는다.
365일 동일 수준을 가정하여
$24\times365=8{,}760$건/대·년으로 환산하고,
기본값의 $\pm25\%$를 민감도 분석 범위로 설정한다.

\begingroup
\small
\setlength{\tabcolsep}{4pt}
\renewcommand{\arraystretch}{1.2}

\begin{longtable}{@{}
    P{0.12\linewidth}
    P{0.18\linewidth}
    P{0.20\linewidth}
    P{0.43\linewidth}
@{}}
\toprule
설정
& 하루 환산 \newline (건/대·일)
& 연간 $\kappa_c$ \newline (건/대·년)
& 설정 근거 \\
\midrule
\endfirsthead

\toprule
설정
& 하루 환산 \newline (건/대·일)
& 연간 $\kappa_c$ \newline (건/대·년)
& 설정 근거 \\
\midrule
\endhead

낮은값
& 18
& 6,570
& 기본값의 $75\%$.
연구자가 정한 민감도 분석 가정. \\
\addlinespace

\textbf{기본값}
& \textbf{24}
& \textbf{8,760}
& Waymo 공개 총량의 단순 비율
약 23.8건/대·일을 반올림하고,
365일을 곱하여 환산. \\
\addlinespace

높은값
& 30
& 10,950
& 기본값의 $125\%$.
연구자가 정한 민감도 분석 가정. \\
\bottomrule
\end{longtable}

\endgroup

%%%%%%%%%%%%%%%%%%%%%%%%%%%%%%%%%%%%

\subsection*{\hypertarget{rationale-R3}{R3. 연간 호출서비스 이용 건수: SFCTA 보고자료}}
\noindent
\textbf{O:}
SFCTA는 2019년 9월부터 2020년 8월까지 12개월 동안
San Francisco 카운티의 Uber 및 Lyft 통행 건수를
38,527,207건으로 집계하였다.
이는 보고서의 Figure 5에 사용된 공개 집계값이다.\\
\textbf{적용:}
San Francisco의 연간 호출서비스 이용 규모를 나타내는
참고값으로 사용하고, R4의 교통수단 분담률과 결합하여
연간 잠재 이용 기회 $M_c$의 실험용 기본값을 설정한다.\\
\textbf{한계:}
코로나19의 영향을 받은 기간이 포함되어 있으며,
보고서에서 사업자의 집계보고와 개별 운행보고 간
불일치를 지적하고 있다.
또한 Uber 및 Lyft의 보고 통행 건수이므로,
현재의 로보택시와 일반 택시까지 포함한
전체 호출서비스 이용 건수로 해석하지 않는다.
\par \href{https://tncs2020.sfcta.org/ch3_general/}
{SFCTA, \textit{TNCs 2020: A Profile of Ride-Hailing
in California}, Chapter 3, Figure 5
(San Francisco 카운티의 12개월 통행 건수)}.

%%%%%%%%%%%%%%%%%%%%%%%%%%%%%%%%%%%%%%%%%%
\subsection*{\hypertarget{rationale-R4}{R4. 교통수단 분담률과 연간 잠재 이용 기회 산출}}
\noindent
\textbf{O:}
SFMTA의 2019년 통행 조사에서 Uber 및 Lyft 등
TNC의 분담률은 전체 조사 통행의 약 5\%로 제시되었다.
다른 교통수단의 분담률은 개인 자동차 48\%,
도보 22\%, 대중교통 등 Transit 22\%,
자전거 2\%, 기타 1\%이며, 반올림된 값이다.
일반 택시는 기타 범주에 포함된다.
이 비율은 호출서비스 이용 경험자의 비율이 아니라
통행 건수를 기준으로 한 교통수단 분담률이다.\\
\textbf{적용:}
TNC 통행 비율을 5\%, 그 외 통행 비율을 95\%로 두고,
R3의 연간 TNC 통행 건수에 대응하는
전체 통행 규모를 역산한다.
첫 실험에서는 이렇게 역산한 전체 통행 규모를
연간 잠재 이용 기회 $M_c$의 대리값으로 사용한다.\\
\textbf{한계:}
이 조사는 Bay Area 주민의 San Francisco 내부 및
San Francisco를 오가는 통행을 대상으로 하며,
Bay Area 외부 방문객과 상업 통행은 제외한다.
따라서 R3의 사업자 보고자료와 조사 시기,
모집단 및 집계 단위가 완전히 일치하지 않는다.
특히 사업자의 운행 건수와 개인의 통행 건수 사이에는
동승 등에 따른 차이가 있을 수 있다.
두 자료를 결합한 결과는 이러한 차이를 허용한
실험용 추정값이며, 잠재 이용 기회를 직접 관측한 값이 아니다.
실현되지 않은 이동 기회도 이 계산에 포함되지 않는다.
\par \href{https://www.sfmta.com/media/20810/download?inline=}
{SFMTA, \textit{Travel Decisions Survey 2019:
Summary Report}, pp.~3--5
(조사 범위 및 교통수단 분담률)}. 
\medskip \\
\textbf{D: 산출 및 실험 설정.}
연간 호출서비스 통행 건수에 더할 비호출 통행 규모는
단순히 95\%를 곱한 값이 아니라,
비호출 통행과 호출 통행의 비율인 $95/5$를 곱한 값이다.
따라서 기본값은 다음과 같이 설정한다.
\[
\begin{aligned}
M_c
&=38{,}527{,}207
  +38{,}527{,}207\times\frac{0.95}{0.05}
=\frac{38{,}527{,}207}{0.05} =770{,}544{,}140
\quad \text{건/년}.
\end{aligned}
\]
이 계산에서 사용한 TNC 분담률 5\%는
$M_c$를 역산하기 위한 참고 비율이며,
우리 회사의 로보택시 선택확률 $p_c^s$와는 다르다.
기본값의 $\pm25\%$를 민감도 분석 범위로 설정한다.
이 범위는 통계적 신뢰구간이 아니라 연구자가 정한
시장규모 변동 가정이다.

\begingroup
\small
\setlength{\tabcolsep}{4pt}
\renewcommand{\arraystretch}{1.2}
\begin{longtable}{@{}
P{0.12\linewidth}
P{0.18\linewidth}
P{0.24\linewidth}
P{0.39\linewidth}
@{}}
\toprule
설정
& 기본값 대비
& 연간 $M_c$ \newline (건/년)
& 설정 근거 \\
\midrule
\endfirsthead
\toprule
설정
& 기본값 대비
& 연간 $M_c$ \newline (건/년)
& 설정 근거 \\
\midrule
\endhead
낮은값
& 75\%
& 577,908,105
& 기본값의 $75\%$.
연구자가 정한 민감도 분석 가정. \\
\addlinespace
\textbf{기본값}
& \textbf{100\%}
& \textbf{770,544,140}
& SFCTA 연간 보고량 38,527,207건을
SFMTA의 TNC 분담률 5\%로 나누어 역산. \\
\addlinespace
높은값
& 125\%
& 963,180,175
& 기본값의 $125\%$.
연구자가 정한 민감도 분석 가정. \\
\bottomrule
\end{longtable}
\endgroup

\noindent
민감도 분석에서는 다른 입력값과
$p_c^s$의 시나리오를 동일하게 유지하고,
$M_c$만 위의 세 수준으로 변경하여
도시 진입, 차량 배치, 미충족 수요 및
목적함수 값의 변화를 비교한다.

%%%%%%%%%%%%%%%%%%%%%%%%%%%%%%%%%%%%%%%%%%%%%%%%%
\subsection*{\hypertarget{rationale-R5}{R5. 도시 진입 초기비용: Waymo Austin 시설 구축 사례}}
\noindent
\textbf{E:}
텍사스 주 TDLR의 프로젝트 등록자료는 Waymo Austin의
Phase 1 공사 예정비를 \$300,000,
Phase 2 공사 예정비를 \$2,500,000으로 제시한다.
Phase 1은 외부 EV 충전기 3개와 관련 전기공사이며,
Phase 2는 기존 창고 내 사무실 및 사업 운영 공간의 개조공사로
약 38,000~ft$^2$를 대상으로 한다.
두 프로젝트는 같은 주소에 2023년 12월 22일 등록되었으며,
공사 일정은 2024년으로 제시되어 있다.
이 금액은 실제 최종 지출액이 아니라 등록된 공사 예정비이다. \\
\textbf{적용:}
기존 건물을 임차하여 운영 거점 하나를 개조하는 진입 방식을 가정하고,
두 단계의 공사 예정비 합계를 $F_c$ 산정을 위한 시설 비용 기준액으로 사용한다.
$F_c$는 도시 $c$에 신규 진입할 때 이번 1년 평가기간에
한 번 지불하는 추가 초기비용이며, 연간등가비용이 아니다. \\
\textbf{한계:}
Austin의 단일 시설 사례이므로 다른 도시의 실제 진입비용을 직접 나타내지 않는다.
차량 구매비, 부지 매입비, 연중 임대료 및 반복 운영비를 포함한 계산이 아니며,
해당 시설과 충전기 3개가 배치 차량 전체를 지원할 수 있는지도 검증되지 않았다.
공사 예정비의 실제 지급 시점은 확인되지 않았으며,
실험에서는 이를 진입 첫해에 지불한다고 가정한다.
과거 등록금액은 물가 조정 없이 규모 참고값으로 사용한다.
\medskip
\noindent
\par\href{https://www.tdlr.texas.gov/TABS/Search/Print/TABS2024008167}
{Texas Department of Licensing and Regulation,
\textit{Waymo Austin--Phase 1}, TABS2024008167,
Estimated Cost 및 Scope of Work};
\href{https://www.tdlr.texas.gov/TABS/Search/Print/TABS2024008168}
{같은 기관, \textit{Waymo Austin--Phase 2}, TABS2024008168,
Estimated Cost 및 Scope of Work}.
\medskip \\

\textbf{D/A: 산출 및 실험 설정.}
범위가 구분된 두 공사 단계의 예정비를 합산하여
시설 비용 기준액을 \$2,800,000으로 둔다.
공개자료에서 금액을 확인하지 못한 허가 대응, 교육 및 개시 전 준비 등
추가 일회성 비용에 대해서는 기준액의 25\%를 포괄 가산한다고 가정한다.
이 가산율은 출처의 추정치나 업계 표준이 아닌 연구자 가정이며,
개별 준비비 항목을 조사하여 합산한 값은 아니다.
\[
\begin{aligned}
\text{시설 비용 기준액}
&=300{,}000+2{,}500{,}000=2{,}800{,}000\quad\text{USD},\\
\text{추가 준비비 가정}
&=2{,}800{,}000\times0.25=700{,}000\quad\text{USD},\\
F_c&=2{,}800{,}000+700{,}000
=3{,}500{,}000\quad\text{USD}.
\end{aligned}
\]

\begingroup
\small
\setlength{\tabcolsep}{4pt}
\renewcommand{\arraystretch}{1.2}
\begin{longtable}{@{}
P{0.12\linewidth}
P{0.18\linewidth}
P{0.24\linewidth}
P{0.39\linewidth}
@{}}
\toprule
설정 & 기본값 대비 & 초기 $F_c$ \newline (USD/도시 진입 1회) & 설정 근거 \\
\midrule
\endfirsthead
\toprule
설정 & 기본값 대비 & 초기 $F_c$ \newline (USD/도시 진입 1회) & 설정 근거 \\
\midrule
\endhead
낮은값 & 75\% & \$2,625,000
& 기본값의 $75\%$. 연구자가 정한 민감도 분석 가정. \\
\addlinespace
\textbf{기본값} & \textbf{100\%} & \textbf{\$3,500,000}
& Waymo Austin 공사 예정비 합계 \$2,800,000에
추가 준비비 25\%를 가정하여 가산. \\
\addlinespace
높은값 & 125\% & \$4,375,000
& 기본값의 $125\%$. 연구자가 정한 민감도 분석 가정. \\
\bottomrule
\end{longtable}
\endgroup
\noindent

%%%%%%%%%%%%%%%%%%%%%%%%%%%%%%%%%%%%%%%%
\subsection*{\hypertarget{rationale-R6}{R6. 차량당 연간 비용: AAA (2025년) 보유비용 참고}}
\noindent
\textbf{E:} 미국의 자동차 이용자 단체인 AAA (American Automobile Association)의 개인용 신차 비용 보고서는 연간 감가상각비 \$4,334, 금융비용 \$1,131, 보험료 \$1,694, 면허·등록·세금 \$813을 제시한다. 9개 차종 범주에서 각각 인기 모델 5개를 대상으로 하며, 5년/75,000마일 보유 기준으로 계산한 평균 비용이다.\\
\textbf{적용:} 위 네 항목의 합계를 차량 1대의 준비·배치 및 1년간 추가 고정 보유비용 $K_c$의 합성 실험 기본값으로 채택한다.\\
\textbf{한계:} 개인용 차량 자료이며 로보택시의 실제 비용이나 엄밀한 하한이 아니다. 준비·배치비, 자율주행 장비, 영업용 보험 차이, 고정 주차·관리비를 별도로 보정하거나 항목별로 합산한 것은 아니다. 통행량에 따라 달라지는 연료·정비·타이어 비용은 이 합계에서 제외한다.
\par\href{https://newsroom.aaa.com/wp-content/uploads/2025/09/UPDATE-AAA-Fact-Sheet-Your-Driving-Cost-9.2025-1.pdf#page=1}{AAA, \textit{Your Driving Costs 2025}, 1쪽 비용 항목 표 및 2쪽 Methodology.}

\medskip
\noindent
\textbf{D/A: 산출 및 실험 설정.}
연간 비용 합계는 $4{,}334+1{,}131+1{,}694+813=7{,}972$달러/대·년이다.
첫 합성 실험에서는 이 참고금액을 $K_c$ 전체의 근사값으로 가정한다.
이미 연간 금액이므로 365를 다시 곱하지 않는다.
민감도 분석에서는 기본값의 $\pm25\%$를 적용하여
$7{,}972\times0.75=5{,}979$와 $7{,}972\times1.25=9{,}965$를 비교한다.

\begingroup
\small
\setlength{\tabcolsep}{4pt}
\renewcommand{\arraystretch}{1.2}
\begin{longtable}{@{}P{0.12\linewidth}P{0.18\linewidth}P{0.24\linewidth}P{0.39\linewidth}@{}}
\toprule
설정 & 기본값 대비 & 연간 $K_c$ \newline (USD/대·년) & 설정 근거 \\
\midrule
\endfirsthead
\toprule
설정 & 기본값 대비 & 연간 $K_c$ \newline (USD/대·년) & 설정 근거 \\
\midrule
\endhead
낮은값 & 75\% & \$5,979 & 기본값의 $75\%$. 연구자가 정한 민감도 분석 가정. \\
\addlinespace
\textbf{기본값} & \textbf{100\%} & \textbf{\$7,972} & AAA의 감가상각·금융·보험·면허·등록·세금 비용 합계를 근사값으로 채택. \\
\addlinespace
높은값 & 125\% & \$9,965 & 기본값의 $125\%$. 연구자가 정한 민감도 분석 가정. \\
\bottomrule
\end{longtable}
\endgroup

%%%%%%%%%%%%%%%%%%%%%%%%%%%%%%%%%%%%%%
\subsection*{\hypertarget{rationale-R7}{R7. 운임 비교: Obi의 원본 가격 조사 (2025년 6월)}}
\textbf{O:} 2025년 3월 25일--4월 25일 SF 검색 중 세 제공업체를 모두 이용할 수 있었던 705개의 표준화된 비교에서 Waymo의 평균 제시 운임은 \$20.43이었다. \textbf{적용:} 운임 규모의 참고값이다. \textbf{한계:} 제시 운임은 실제 기업 수입이나 공헌이익이 아니다. 플랫폼 배분 비율과 추가 비용은 여전히 알 수 없다. \textbf{A의 예:} 기업 귀속 운임 \$18에서 추가 비용 \$8을 빼면 $c_c^-=\$10$이다. 이 계산은 예시이며 Obi에서 추정한 값이 아니다.
\par\href{https://rideobi.com/wp-content/uploads/2025/10/Obi-Report-Pricing-Insights-On-Waymo-Uber-and-Lyft.pdf}{Obi: The Road Ahead, 요약 및 문서에 표시된 5쪽.}

\section{실행한 실험과의 연결}
기존 일 단위 실험에서 실제로 선택한 입력값, 시나리오 정의, 실행 절차 및 결과의 기준 문서는 \texttt{numerical\_experiments.tex}이다. 현재 Overleaf에서는 \texttt{1. main(modeling)/baseline model/} 폴더에 있다. 위 표의 연간 입력값은 제안 단계이며 아직 실행한 연 단위 실험의 결과가 아니다. 이후 실험 파일과 코드를 연 단위 입력값으로 바꾼 뒤 실행 식별자와 결과를 기록해야 한다.

일반적인 일관성 점검으로 총 가용 차량 수와 $\sum_c M_c\bar p_c/\kappa_c$를 비교하고, 미진입 비용을 수식으로 계산하며, 추가 서비스의 공헌이익이 배치비용을 충당할 수 있는지 확인한다. 이 점검은 실행 전에 수행해야 하며 최적화 결과를 대신하지 않는다.
진입이 경제적으로 의미 있으려면 감소하는 미충족 수요 비용이 고정비용과 차량비용을 충당할 수 있어야 한다. $c_c^-\kappa_c>K_c$이면 충분한 수요를 처리하는 차량은 진입비용을 고려하기 전 양의 운영 공헌이익을 낼 수 있다. 그렇다고 차량 수 제약 아래에서 진입이 최적이라는 사실까지 입증되는 것은 아니다. 도시별로 정확히 점검하는 식은 다음과 같다.
\[
\max_{1\le q\le\min(\bar q_c,Q)}\{g_c(0)-g_c(q)\}>0,
\]
\[
g_c(q)=F_c\mathbf1_{q>0}+K_cq+
\frac{1}{|S|}\sum_{s\in S}[c_c^+(\kappa_cq-D_c^s)^++c_c^-(D_c^s-\kappa_cq)^+].
\]
도시들은 가용 차량 $Q$를 두고 경쟁하므로 도시별 수익성을 따로 확인하는 것만으로 전체 최적해가 결정되지는 않는다.

\section{해석의 범위}
진입하지 않은 도시의 잠재 수요도 미충족 수요로 남는다. $y_c=0$이면 $q_c=0$이고 $U_c^s=D_c^s$이다. $c_c^-$를 상실한 공헌이익으로 해석하면 미진입 비용은 상수인 수익 기회의 비교 기준이다. 이 상수를 목적함수에서 빼면 모델이 기대 서비스 공헌이익에서 고정비용, 차량비용 및 추가 유휴 비용을 뺀 값을 최대화하는 것과 같음을 알 수 있다. 법적인 서비스 제공 의무를 뜻하지 않는다.

수요는 외생적이고 시나리오 비용은 도시별로 합산되므로, 각 도시의 주변 수요분포와 발생 확률을 유지한 채 도시 간 의존관계만 바꾸면 어떤 고정 배치에 대해서도 기본 목적함수는 변하지 않는다. 따라서 의존관계 실험은 모델 구조의 점검에 해당한다. 재배치나 공동 사후대응처럼 시나리오 내 도시들을 연결하는 의사결정이 추가되어야 상관관계가 결과에 영향을 줄 수 있다.

\textbf{실제 자료에 맞춘 설정을 위해 추가로 필요한 근거:} 대상 서비스 구역과 시점, 기업 귀속 운임과 통행당 추가 비용, 차량 구매·리스 및 영업용 보험 조건, 실제 서비스 시간과 전체 통행 처리 주기, 차고지·충전 제약, 그리고 $M_c$와 분모가 일치하는 선택모형이다. 이를 확보하기 전까지 결과가 보여 주는 것은 모델의 정확성과 합성 조건에서의 비교 행동이며, 실제 시장의 최적 차량 수에 대한 권고는 아니다.
\end{document}
